\section{Yöntem}\label{sec:yontem}

Bütün stratejiler aynı ILS motorunu paylaşır ve yalnız double-bridge vuruşunun kesim konumlarını
nasıl seçtiklerinde ayrılır. Hiçbir model kabul kararı vermez: sinyaller yalnız vuruşun nereye
yapılacağını önerir, kabul her zaman gerçek tur maliyetiyle yapılır. Böylece stratejiler
arasındaki her fark doğrudan hedef seçimine bağlanabilir.

\subsection{Problem gösterimi}\label{sec:yontem-problem}

Bir TSP örneği $I=(V,d)$ ile gösterilsin; $n=|V|$ şehir içerir ve $d$ örneğin bildirdiği uzaklık
işlevidir (EUC\_2D'de yuvarlatılmış Öklid uzaklığı, GEO'da TSPLIB büyük çember uzaklığı). Tur
$T=(v_1,\dots,v_n)$, $v_{n+1}=v_1$ olmak üzere döngüsel bir permütasyondur ve maliyeti
$C(T)=\sum_{i=1}^{n} d(v_i,v_{i+1})$'dir. Kalite, bilinen optimuma ya da referans çözüme (sentetik problemlerde LKH-3'ün 60~sn'lik
koşusu) göre yüzde farkla ölçülür: $\mathrm{gap}(T)=100\,(C(T)-C^\star)/C^\star$. İki strateji arasındaki fark gap
\emph{puanı} olarak verilir; negatif değer soldaki stratejinin lehinedir. $N_k(u)$, $u$'nun $k$
en yakın komşusunun listesidir.

\subsection{Hızlı ILS test ortamı}\label{sec:yontem-hizli}

Bir vuruş politikasının etkisi, ancak vuruş döngüsü bütçe içinde binlerce kez dönebiliyorsa
ölçülebilir. Bu amaçla standart hızlı yerel arama tekniklerini \citep{bentley_1992_fast,
johnson_1997_case} izleyen ve Numba ile makine koduna derlenen bir ILS kurduk (\path{fast_ils.py}).
Aynı dosya Aşama~2 ve Aşama~3'te değiştirilmeden kullanıldı; SHA-256 özeti iki protokolde de
kayıtlıdır.

\begin{itemize}
  \item \textbf{Yerel arama.} 2-opt ve Or-opt (1--3 şehirlik segmentin iki yönde yeniden
  yerleştirilmesi); her şehir için $N_{10}$ aday listesi ve $d(a,c)<d(a,\mathrm{succ}(a))$ budama
  kuralı. 2-opt ters çevirmesi kısa olan yay üzerinde yapılır.
  \item \textbf{Bakma bitleri.} Etkin şehirler bir kuyrukta tutulur. Vuruştan sonra yalnız değişen
  kenarların uçları kuyruğa girer; bir vuruşun bedeli tüm tur taraması yerine yerel onarımla
  sınırlı kalır \citep{bentley_1992_fast}.
  \item \textbf{Kabul.} Aday yalnız $C(T')<C(T)$ ise kabul edilir; aksi hâlde tur geri yüklenir.
  Maliyet artımlı izlenir ve koşu sonunda tam yeniden hesapla doğrulanır.
  \item \textbf{Başlangıç.} En yakın komşu turu (başlangıç şehri tohuma bağlı) ve tam 2-opt/Or-opt.
  Bir (problem, tohum) çiftinde bütün stratejiler aynı başlangıç turunu paylaşır; başlangıç süresi
  bütçeye dahil değildir.
\end{itemize}

Amaç LKH-3 ile yarışmak değil, vuruş politikasının gerçekten devrede olduğu, açık ve yeniden
üretilebilir bir ölçüm düzeneği sağlamaktır.

\subsection{Vuruş geometrisi ve havuz kuralı}\label{sec:yontem-kick}

Double-bridge \citep{martin_1992_largestep, lourenco_2010_ils}, kesim konumları $a<b<c$ için turu
$T_{[:a]}\,T_{[b:c]}\,T_{[a:b]}\,T_{[c:]}$ biçiminde yeniden bağlar. \emph{Global} geometride
kesimler $[1,n-1]$ aralığında, birbirinden en az iki konum uzakta seçilir. Aşama~2 ayrıca
\emph{yerel} (segment) geometriyi de sınadı: ilk kesimden sonra diğer ikisi en fazla $W=50$ konum
içinde seçilir \citep{applegate_2003_clk}. Aşama~3 yalnız global geometriyi kullanır.

Hedefli stratejiler her tazelemede bir \emph{havuz} kurar: sinyale göre en şüpheli
\[
  m=\max\bigl(12,\lfloor f\,n\rfloor\bigr)
\]
tur kenarı. Global vuruşun üç kesiminin her biri, $p$ olasılıkla havuzdan tekdüze seçilen bir
kenarın üzerine konur, yani o kenar kırılır; aksi hâlde kesim tekdüzedir. Havuz oranı $f$ ve hedef
olasılığı $p$ havuz kuralını tanımlar. Varsayılan kural $f=0{,}06$, $p=0{,}75$'tir; Aşama~3'te üç
sinyal için kural geliştirme setlerinde seçilmiştir (Bölüm~\ref{sec:deney-asama3}). Tekdüze
kesimler, hedefli stratejilerde de keşfin sürmesini sağlar.

\subsection{Hedefleme sinyalleri}\label{sec:yontem-sinyal}

Sekiz sinyal üç gruba ayrılır: eğitimsiz geometrik vekiller, çevrimdışı eğitilmiş skorlar ve koşu
içinde öğrenen skor. Hepsi aynı havuz kuralını ve aynı maliyet muhasebesini kullanır.

\paragraph{Rastgele.} Havuz yoktur; bütün kesimler tekdüzedir. Bu, ILS'nin klasik davranışıdır ve
ana kontrol koludur.

\paragraph{kNN sırası.} Tur kenarı $(u,v)$ için şüphe skoru
$s_{\mathrm{kNN}}(u,v)=\tfrac12\bigl(r_u(v)+r_v(u)\bigr)$'dir; $r_u(v)$, $v$'nin $N_{24}(u)$
içindeki sırasıdır (listede yoksa $30$). Eşitlikler, kenar uzunluğunun uç noktaların en yakın komşu
uzaklıklarına oranıyla bozulur. Skor hiçbir eğitim ya da çözücü etiketi kullanmaz.

\paragraph{Fazla-uzunluk vekili.} İyi bir turda bir şehre bağlanan kenarlar tipik olarak en yakın
komşularına giden kenarlar ölçeğindedir. Bu ölçekten ne kadar uzun olduğu, bir kenarı kırmanın
olası kazancı için kaba ama ucuz bir göstergedir. Vekil
\begin{equation}
  e(u,v)=d(u,v)-\tfrac12\bigl(d_{\mathrm{ref}}(u)+d_{\mathrm{ref}}(v)\bigr),\qquad
  d_{\mathrm{ref}}(u)=\tfrac12\bigl(d(u,u_{(1)})+d(u,u_{(2)})\bigr)
  \label{eq:excess}
\end{equation}
ile tanımlanır; $u_{(1)}$ ve $u_{(2)}$, $u$'nun birinci ve ikinci en yakın komşusudur. Havuz,
$e$ değeri en büyük $m$ kenardır. $d_{\mathrm{ref}}$ örnek başına bir kez hesaplanır; böylece
tazeleme kenar başına $O(1)$'dir. Vekil, aşağıdaki kazanç hedefinin eğitimsiz karşılığı olarak
tasarlandı ve GLS'nin uzun kenar ilkesini \citep{voudouris_1999_gls} vuruş hedeflemeye taşır. Sıra
tabanlı kNN skorundan farkı, kenarın komşu listesindeki yerine değil, uzunluğun büyüklüğüne
bakmasıdır. Protokol sonrasında, keşifsel olarak, $d_{\mathrm{ref}}$ terimi atılmış ham uzunluk
sinyali ($e(u,v)=d(u,v)$; GLS'nin cezasız fayda ölçütü) de sınandı. Bu kol, fazla-uzunluk için
geliştirmede seçilen havuz kuralını (0,06/0,5) olduğu gibi kullandı. % E751

\paragraph{GNN.} Örnekle ve topla (GraphSAGE) tarzında, saf NumPy ile gerçeklenmiş küçük bir
ileti geçişli ağdır: düğüm başına 10 öznitelik (komşu tur kenarlarının uzunlukları, kaldırma
kazancı, en yakın komşu uzaklıkları, dönüş açısı, segmentten sapma ve tur komşularının kNN
sırası), tur kenarı başına 2 öznitelik (z-normalize uzunluk ve komşu sırası), 3 ileti geçişi
katmanı, 16 boyutlu gizli durum ve $N_{24}$ üzerinde geometrik ileti geçişi. Kenar çıkışı, bir tur
kenarının LKH-3 \citep{helsgaun_2017_lkh3} turunda bulunma olasılığını $\hat p_e$ olarak tahmin eder.
Model $n\in[101,400]$ aralığındaki 120 sentetik problemde (beş dağılım) 954 tur anlık görüntüsüyle
eğitilmiş ve dondurulmuştur. Havuz, $\hat p_e$ değeri en düşük $m$ kenardır. Modelin girdileri
arasında kNN sırası doğrudan bulunduğundan, öğrenmenin bu sinyalin ötesine geçip geçmediği ampirik
bir sorudur.

Bu GNN yayımlanmış bir modelin yeniden üretimi değil, yazarın kendi gerçeklemesidir; LKH turlarıyla
eğitilmiş kenar skorları sınıfının küçük bir temsilcisi olarak kullanılır. İşlevselliği iki ölçümle
gösterilir: yanlış kenarları kNN'den daha iyi ayırt eder (Bölüm~\ref{sec:bulgular-asama2}) ve iki
aşamada da rastgele vuruştan iyidir. % E889, E803, E611
Test problemleri ($n=1.376$--$9.882$) eğitim aralığının çok üstündedir; bu nedenle karşılaştırma
bir boyut genellemesi sınamasını da içerir ve vardığımız sonuçlar bu ölçekteki öğrenilmiş skorlarla
sınırlıdır. % E797, E598

\paragraph{Üyelik-MLP ve kazanç-MLP.} GNN'nin kaybının öğrenmeden mi maliyetten mi geldiğini
ayırmak için, GNN'nin eğitim verisinde iki küçük MLP (17 girdi, 32 ve 16 gizli birim) eğitildi.
Girdiler $N_{24}$ listelerinden kenar başına $O(1)$'de hesaplanır: iki uçtaki karşılıklı sıralar ve
bunların en küçük/en büyük değeri, kenar uzunluğunun 1., 5. ve 24. komşu uzaklığına oranları,
komşu iki tur kenarının uzunluk oranları ve sıraları. \emph{Üyelik-MLP} kenarın referans turda
bulunmamasını sınıflandırır. \emph{Kazanç-MLP} ise
\[
  y_e=\mathbb{1}[e\notin T^\star]\cdot\max\bigl(0,\,d_e-\bar d^{\,\star}_e\bigr)/\tilde d
\]
hedefini $\log(1+y_e)$ üzerinden regresyonla öğrenir; $\bar d^{\,\star}_e$, $e$'nin iki ucundaki
referans tur kenarlarının ortalama uzunluğu, $\tilde d$ medyan tur kenarı uzunluğudur. Kazanç
hedefi, vuruşun yanlış kenarı kırmasının yanında bunun ne kadar kazandırabileceğini de ölçer;
Denklem~\eqref{eq:excess}, $\bar d^{\,\star}$ yerine eğitimsiz $d_{\mathrm{ref}}$ koyarak bu hedefe
yaklaşır.

\paragraph{Sabit doğrusal ön-bilgi.} Aynı 17 öznitelik ve bir sabit terim üzerinde, kazanç hedefine
ridge regresyonuyla uydurulmuş doğrusal bir skor $w_0^\top x_e$'dir. Koşu boyunca değişmez; havuz en
yüksek skorlu $m$ kenardır. Bu kol, çevrimiçi öğrenmenin etkisini modelden ayırmak için eşleşik
kontroldür.

\paragraph{Çevrimiçi (Thompson-ridge).} Aynı doğrusal model koşu sırasında, her vuruşun gerçekleşen
sonucundan öğrenir. Kesimleri $e_1,e_2,e_3$ olan bir vuruşun ödülü, yerel arama sonrasındaki
maliyet değişiminden
\begin{equation}
  r=-\frac{C(T_{\text{sonra}})-C(T_{\text{önce}})}{\tilde d_0}
  \label{eq:reward}
\end{equation}
olarak hesaplanır; $\tilde d_0$ başlangıç turunun medyan kenar uzunluğudur. İyileşme pozitif ödül
verir. Ödül toplamsal modellenir, $r\approx w^\top(x_{e_1}+x_{e_2}+x_{e_3})$, ve $w$ çevrimiçi
ridge regresyonuyla güncellenir: $A=\lambda I+\sum xx^\top$, $b=\lambda w_0+\sum r\,x$,
$\lambda=50$. Her tazelemede $w\sim\mathcal N(A^{-1}b,\ \nu^2A^{-1})$, $\nu=0{,}5$, örneklenir ve
tahmini ödülü en yüksek $m$ kenar havuzu oluşturur. Model $w_0$ ile ılık başlatılır. Havuz, diğer sinyallerdeki genel kuralla
(Bölüm~\ref{sec:yontem-maliyet}) her kabul edilen vuruşta, ayrıca 2.000 vuruşta bir yeniden kurulur. İşaretli ödülün bilinçli bir gerekçesi
vardır: vuruşların çok küçük bir kısmı kabul edildiğinden, $\max(0,\cdot)$ biçimindeki bir ödül
neredeyse her vuruş için sıfır olurdu. Ödül işareti ve havuz seçimi birim testleriyle doğrulandı:
iyileştiren vuruş pozitif, nötr vuruş sıfır ödül alır ve havuz en yüksek tahmini ödülü seçer. % E959
Aşama~3 geliştirme koşularında soğuk başlatılan ($w_0=0$) bir sürüm de sınandı.

\subsection{Maliyet muhasebesi}\label{sec:yontem-maliyet}

Havuz, incumbent tur değiştiğinde yeniden hesaplanır. Skorlama süresi (GNN çıkarımı, MLP ileri
geçişi, vekil hesabı, Thompson örneklemesi) duvar-saati bütçesinden düşülür; yani her sinyal
ürettiği bilginin bedelini aynı bütçeden öder. Toplam skorlama süresi geçen sürenin \%10'unu
aşarsa tazeleme ertelenir. Bu sırada havuzdaki bir kenar artık turda değilse o kesim tekdüze
seçilir; buna \emph{bayat kesim} diyoruz. Örnek başına bir kez yapılan hazırlıklar ($N_{24}$
listeleri, komşu uzaklıkları, ödül ölçeği) hiçbir sinyal için bütçeye yüklenmez. Numba derleme
süresi her işçi sürecinde bir ısınma çağrısıyla bütçe dışına alınır.

\subsection{Aşama 1 iskeleti}\label{sec:yontem-crine}

Aşama~1, öğrenilmiş vuruş hedeflemesini CRINE-ILS (Counterfactual Residual Improvement through Neighborhood Expansion) adlı yavaş bir iskelette sınamıştı: saf Python
yerel arama, her vuruştan sonra tüm tur üzerinde sınırlı onarım ve ikinci bir GNN ile sıralanan
şehirler çevresinde komşuluk genişletme (\path{adaptive_ils.py}). Bu iskelette bir vuruş yüzlerce
milisaniye sürüyordu; kısa bütçede vuruş döngüsü neredeyse hiç çalışmadığı için stratejiler
arasındaki farklar tohum gürültüsüne indirgendi (Bölüm~\ref{sec:bulgular}). Hızlı test ortamı bu
dersin doğrudan sonucudur.

\subsection{İşlem karmaşıklığı}\label{sec:yontem-karmasiklik}

Bir tazelemede fazla-uzunluk vekili, kNN skoru ve iki MLP kenar başına sabit iş yapar; toplam
$O(n)$'dir. GNN çıkarımı da $k=24$ ve gizli boyut sabit olduğundan $O(n)$'dir, ancak sabit çarpanı
birkaç ileti geçişi katmanı nedeniyle büyüktür. Çevrimiçi skorun her vuruştaki güncellemesi 18
boyutlu bir dış çarpım, her tazelemedeki örneklemesi $18\times18$ bir ters alma gerektirir. Global
vuruş ve reddedilen vuruştan sonra tur geri yüklemesi $O(n)$ dizi kopyasıdır; vuruş sonrası onarım,
etkinleşen şehir sayısıyla orantılıdır.
